% figA1_crosstask.tex — Figure A1: Cross-task performance + TDI slopes
% Canonical path for the paper: % figA1_crosstask.tex — Figure A1: Cross-task performance + TDI slopes
% Canonical path for the paper: % figA1_crosstask.tex — Figure A1: Cross-task performance + TDI slopes
% Canonical path for the paper: % figA1_crosstask.tex — Figure A1: Cross-task performance + TDI slopes
% Canonical path for the paper: \input{figures/figA1_crosstask}
%
% Right panel uses internal symbolic y keys (no punctuation in tokens) so
% pgfplots never parses '/' as division; human-readable names come from
% yticklabels. Multiple xbar series keep bar shift auto disabled so bars align
% with ticks.

\begin{figure*}[t]
\centering
\begin{subfigure}[b]{0.52\textwidth}
\begin{tikzpicture}
\begin{axis}[
  pmh bar,
  symbolic x coords={T01,T02,T03,T04,T05,T06,T07},
  xtick=data,
  xticklabel style={font=\footnotesize},
  ylabel={Normalised robust metric (best$=$100)},
  ymin=0, ymax=118,
  legend style={at={(0.5,1.02)}, anchor=south, legend columns=3},
  width=0.92\linewidth, height=5.5cm,
  bar width=4pt,
]
\addplot[fill=colB0, draw=none]
  coordinates {(T01,50)(T02,95)(T03,93)(T04,72)(T05,100)(T06,65)(T07,76)};
\addlegendentry{B0 (ERM)}
\addplot[fill=colVAT, draw=none]
  coordinates {(T01,81)(T02,86)(T03,82)(T04,93)(T05,16)(T06,100)(T07,89)};
\addlegendentry{VAT}
\addplot[fill=colPMH, draw=none]
  coordinates {(T01,100)(T02,100)(T03,100)(T04,100)(T05,81)(T06,97)(T07,100)};
\addlegendentry{E1 (PMH)}
\end{axis}
\end{tikzpicture}
\caption{Cross-task robust performance (normalised per task).}
\end{subfigure}
\hfill
\begin{subfigure}[b]{0.46\textwidth}
\begin{tikzpicture}
\begin{axis}[
  pmh hbar,
  symbolic y coords={mPGD4,mPGD2,mPMH,mNoPMH,mVAT,mB0},
  ytick=data,
  yticklabels={%
    {PGD-4/255},%
    {PGD-2/255},%
    {E1 (PMH)},%
    {E1 no PMH},%
    {VAT},%
    {B0},%
  },
  yticklabel style={font=\footnotesize, align=right},
  xlabel={$\Delta$TDI/$\Delta\sigma$ slope (lower $=$ better)},
  xmin=0, xmax=8.0,
  width=0.78\linewidth, height=5.5cm,
  bar width=5pt,
  /pgfplots/bar shift auto/.style={/pgf/bar shift=0pt},
  bar shift=0pt,
  nodes near coords,
  nodes near coords style={font=\tiny},
  enlarge y limits=0.18,
]
\addplot[fill=colPGD4, draw=none] coordinates {(2.94,mPGD4)};
\addplot[fill=colPGD2, draw=none] coordinates {(3.01,mPGD2)};
\addplot[fill=colPMH,  draw=none] coordinates {(3.20,mPMH)};
\addplot[fill=colNoPMH,draw=none] coordinates {(3.40,mNoPMH)};
\addplot[fill=colVAT,  draw=none] coordinates {(5.09,mVAT)};
\addplot[fill=colB0,   draw=none] coordinates {(6.48,mB0)};
\end{axis}
\end{tikzpicture}
\caption{TDI degradation slope per method.}
\end{subfigure}
\caption{\textbf{Cross-task performance and TDI degradation slopes.}
\emph{Left:} Bars normalise each task so the best method among \{B0, VAT, E1\} is 100
(replication, seed~42): PMH leads five of seven; B0 leads Task~05 (mean PCK@0.05)
and VAT leads Task~06 (avg-shift rank-1). E1\_node on Task~03.
\emph{Right:} PMH and E1 no PMH have the lowest TDI slopes (3.20 and 3.40);
B0/VAT degrade most rapidly (6.48 and 5.09).
PGD's controlled slope masks a damaged floor (TDI@0$=$1.336).}
\label{fig:a1}
\end{figure*}

%
% Right panel uses internal symbolic y keys (no punctuation in tokens) so
% pgfplots never parses '/' as division; human-readable names come from
% yticklabels. Multiple xbar series keep bar shift auto disabled so bars align
% with ticks.

\begin{figure*}[t]
\centering
\begin{subfigure}[b]{0.52\textwidth}
\begin{tikzpicture}
\begin{axis}[
  pmh bar,
  symbolic x coords={T01,T02,T03,T04,T05,T06,T07},
  xtick=data,
  xticklabel style={font=\footnotesize},
  ylabel={Normalised robust metric (best$=$100)},
  ymin=0, ymax=118,
  legend style={at={(0.5,1.02)}, anchor=south, legend columns=3},
  width=0.92\linewidth, height=5.5cm,
  bar width=4pt,
]
\addplot[fill=colB0, draw=none]
  coordinates {(T01,50)(T02,95)(T03,93)(T04,72)(T05,100)(T06,65)(T07,76)};
\addlegendentry{B0 (ERM)}
\addplot[fill=colVAT, draw=none]
  coordinates {(T01,81)(T02,86)(T03,82)(T04,93)(T05,16)(T06,100)(T07,89)};
\addlegendentry{VAT}
\addplot[fill=colPMH, draw=none]
  coordinates {(T01,100)(T02,100)(T03,100)(T04,100)(T05,81)(T06,97)(T07,100)};
\addlegendentry{E1 (PMH)}
\end{axis}
\end{tikzpicture}
\caption{Cross-task robust performance (normalised per task).}
\end{subfigure}
\hfill
\begin{subfigure}[b]{0.46\textwidth}
\begin{tikzpicture}
\begin{axis}[
  pmh hbar,
  symbolic y coords={mPGD4,mPGD2,mPMH,mNoPMH,mVAT,mB0},
  ytick=data,
  yticklabels={%
    {PGD-4/255},%
    {PGD-2/255},%
    {E1 (PMH)},%
    {E1 no PMH},%
    {VAT},%
    {B0},%
  },
  yticklabel style={font=\footnotesize, align=right},
  xlabel={$\Delta$TDI/$\Delta\sigma$ slope (lower $=$ better)},
  xmin=0, xmax=8.0,
  width=0.78\linewidth, height=5.5cm,
  bar width=5pt,
  /pgfplots/bar shift auto/.style={/pgf/bar shift=0pt},
  bar shift=0pt,
  nodes near coords,
  nodes near coords style={font=\tiny},
  enlarge y limits=0.18,
]
\addplot[fill=colPGD4, draw=none] coordinates {(2.94,mPGD4)};
\addplot[fill=colPGD2, draw=none] coordinates {(3.01,mPGD2)};
\addplot[fill=colPMH,  draw=none] coordinates {(3.20,mPMH)};
\addplot[fill=colNoPMH,draw=none] coordinates {(3.40,mNoPMH)};
\addplot[fill=colVAT,  draw=none] coordinates {(5.09,mVAT)};
\addplot[fill=colB0,   draw=none] coordinates {(6.48,mB0)};
\end{axis}
\end{tikzpicture}
\caption{TDI degradation slope per method.}
\end{subfigure}
\caption{\textbf{Cross-task performance and TDI degradation slopes.}
\emph{Left:} Bars normalise each task so the best method among \{B0, VAT, E1\} is 100
(replication, seed~42): PMH leads five of seven; B0 leads Task~05 (mean PCK@0.05)
and VAT leads Task~06 (avg-shift rank-1). E1\_node on Task~03.
\emph{Right:} PMH and E1 no PMH have the lowest TDI slopes (3.20 and 3.40);
B0/VAT degrade most rapidly (6.48 and 5.09).
PGD's controlled slope masks a damaged floor (TDI@0$=$1.336).}
\label{fig:a1}
\end{figure*}

%
% Right panel uses internal symbolic y keys (no punctuation in tokens) so
% pgfplots never parses '/' as division; human-readable names come from
% yticklabels. Multiple xbar series keep bar shift auto disabled so bars align
% with ticks.

\begin{figure*}[t]
\centering
\begin{subfigure}[b]{0.52\textwidth}
\begin{tikzpicture}
\begin{axis}[
  pmh bar,
  symbolic x coords={T01,T02,T03,T04,T05,T06,T07},
  xtick=data,
  xticklabel style={font=\footnotesize},
  ylabel={Normalised robust metric (best$=$100)},
  ymin=0, ymax=118,
  legend style={at={(0.5,1.02)}, anchor=south, legend columns=3},
  width=0.92\linewidth, height=5.5cm,
  bar width=4pt,
]
\addplot[fill=colB0, draw=none]
  coordinates {(T01,50)(T02,95)(T03,93)(T04,72)(T05,100)(T06,65)(T07,76)};
\addlegendentry{B0 (ERM)}
\addplot[fill=colVAT, draw=none]
  coordinates {(T01,81)(T02,86)(T03,82)(T04,93)(T05,16)(T06,100)(T07,89)};
\addlegendentry{VAT}
\addplot[fill=colPMH, draw=none]
  coordinates {(T01,100)(T02,100)(T03,100)(T04,100)(T05,81)(T06,97)(T07,100)};
\addlegendentry{E1 (PMH)}
\end{axis}
\end{tikzpicture}
\caption{Cross-task robust performance (normalised per task).}
\end{subfigure}
\hfill
\begin{subfigure}[b]{0.46\textwidth}
\begin{tikzpicture}
\begin{axis}[
  pmh hbar,
  symbolic y coords={mPGD4,mPGD2,mPMH,mNoPMH,mVAT,mB0},
  ytick=data,
  yticklabels={%
    {PGD-4/255},%
    {PGD-2/255},%
    {E1 (PMH)},%
    {E1 no PMH},%
    {VAT},%
    {B0},%
  },
  yticklabel style={font=\footnotesize, align=right},
  xlabel={$\Delta$TDI/$\Delta\sigma$ slope (lower $=$ better)},
  xmin=0, xmax=8.0,
  width=0.78\linewidth, height=5.5cm,
  bar width=5pt,
  /pgfplots/bar shift auto/.style={/pgf/bar shift=0pt},
  bar shift=0pt,
  nodes near coords,
  nodes near coords style={font=\tiny},
  enlarge y limits=0.18,
]
\addplot[fill=colPGD4, draw=none] coordinates {(2.94,mPGD4)};
\addplot[fill=colPGD2, draw=none] coordinates {(3.01,mPGD2)};
\addplot[fill=colPMH,  draw=none] coordinates {(3.20,mPMH)};
\addplot[fill=colNoPMH,draw=none] coordinates {(3.40,mNoPMH)};
\addplot[fill=colVAT,  draw=none] coordinates {(5.09,mVAT)};
\addplot[fill=colB0,   draw=none] coordinates {(6.48,mB0)};
\end{axis}
\end{tikzpicture}
\caption{TDI degradation slope per method.}
\end{subfigure}
\caption{\textbf{Cross-task performance and TDI degradation slopes.}
\emph{Left:} Bars normalise each task so the best method among \{B0, VAT, E1\} is 100
(replication, seed~42): PMH leads five of seven; B0 leads Task~05 (mean PCK@0.05)
and VAT leads Task~06 (avg-shift rank-1). E1\_node on Task~03.
\emph{Right:} PMH and E1 no PMH have the lowest TDI slopes (3.20 and 3.40);
B0/VAT degrade most rapidly (6.48 and 5.09).
PGD's controlled slope masks a damaged floor (TDI@0$=$1.336).}
\label{fig:a1}
\end{figure*}

%
% Right panel uses internal symbolic y keys (no punctuation in tokens) so
% pgfplots never parses '/' as division; human-readable names come from
% yticklabels. Multiple xbar series keep bar shift auto disabled so bars align
% with ticks.

\begin{figure*}[t]
\centering
\begin{subfigure}[b]{0.52\textwidth}
\begin{tikzpicture}
\begin{axis}[
  pmh bar,
  symbolic x coords={T01,T02,T03,T04,T05,T06,T07},
  xtick=data,
  xticklabel style={font=\footnotesize},
  ylabel={Normalised robust metric (best$=$100)},
  ymin=0, ymax=118,
  legend style={at={(0.5,1.02)}, anchor=south, legend columns=3},
  width=0.92\linewidth, height=5.5cm,
  bar width=4pt,
]
\addplot[fill=colB0, draw=none]
  coordinates {(T01,50)(T02,95)(T03,93)(T04,72)(T05,100)(T06,65)(T07,76)};
\addlegendentry{B0 (ERM)}
\addplot[fill=colVAT, draw=none]
  coordinates {(T01,81)(T02,86)(T03,82)(T04,93)(T05,16)(T06,100)(T07,89)};
\addlegendentry{VAT}
\addplot[fill=colPMH, draw=none]
  coordinates {(T01,100)(T02,100)(T03,100)(T04,100)(T05,81)(T06,97)(T07,100)};
\addlegendentry{E1 (PMH)}
\end{axis}
\end{tikzpicture}
\caption{Cross-task robust performance (normalised per task).}
\end{subfigure}
\hfill
\begin{subfigure}[b]{0.46\textwidth}
\begin{tikzpicture}
\begin{axis}[
  pmh hbar,
  symbolic y coords={mPGD4,mPGD2,mPMH,mNoPMH,mVAT,mB0},
  ytick=data,
  yticklabels={%
    {PGD-4/255},%
    {PGD-2/255},%
    {E1 (PMH)},%
    {E1 no PMH},%
    {VAT},%
    {B0},%
  },
  yticklabel style={font=\footnotesize, align=right},
  xlabel={$\Delta$TDI/$\Delta\sigma$ slope (lower $=$ better)},
  xmin=0, xmax=8.0,
  width=0.78\linewidth, height=5.5cm,
  bar width=5pt,
  /pgfplots/bar shift auto/.style={/pgf/bar shift=0pt},
  bar shift=0pt,
  nodes near coords,
  nodes near coords style={font=\tiny},
  enlarge y limits=0.18,
]
\addplot[fill=colPGD4, draw=none] coordinates {(2.94,mPGD4)};
\addplot[fill=colPGD2, draw=none] coordinates {(3.01,mPGD2)};
\addplot[fill=colPMH,  draw=none] coordinates {(3.20,mPMH)};
\addplot[fill=colNoPMH,draw=none] coordinates {(3.40,mNoPMH)};
\addplot[fill=colVAT,  draw=none] coordinates {(5.09,mVAT)};
\addplot[fill=colB0,   draw=none] coordinates {(6.48,mB0)};
\end{axis}
\end{tikzpicture}
\caption{TDI degradation slope per method.}
\end{subfigure}
\caption{\textbf{Cross-task performance and TDI degradation slopes.}
\emph{Left:} Bars normalise each task so the best method among \{B0, VAT, E1\} is 100
(replication, seed~42): PMH leads five of seven; B0 leads Task~05 (mean PCK@0.05)
and VAT leads Task~06 (avg-shift rank-1). E1\_node on Task~03.
\emph{Right:} PMH and E1 no PMH have the lowest TDI slopes (3.20 and 3.40);
B0/VAT degrade most rapidly (6.48 and 5.09).
PGD's controlled slope masks a damaged floor (TDI@0$=$1.336).}
\label{fig:a1}
\end{figure*}
